\documentclass[12pt,a4paper]{article}

\usepackage[utf8]{inputenc}
\usepackage[T1]{fontenc}
\usepackage[french]{babel}
\usepackage[margin=2.5cm]{geometry}
\usepackage{listings}
\usepackage{xcolor}
\usepackage{hyperref}
\usepackage{parskip}
\usepackage{booktabs}
\usepackage{fancyhdr}
\usepackage{lmodern}

\setlength{\headheight}{14.5pt}

% ── style des blocs de code ───────────────────────────────────────────────────
\definecolor{codebg}{HTML}{F5F5F5}
\definecolor{keyword}{HTML}{0000CC}
\definecolor{comment}{HTML}{007700}
\definecolor{string}{HTML}{CC0000}
\definecolor{linenum}{HTML}{888888}

\lstdefinestyle{cstyle}{
    language=C,
    backgroundcolor=\color{codebg},
    basicstyle=\ttfamily\small,
    keywordstyle=\color{keyword}\bfseries,
    commentstyle=\color{comment}\itshape,
    stringstyle=\color{string},
    numberstyle=\tiny\color{linenum},
    numbers=left,
    numbersep=8pt,
    stepnumber=1,
    breaklines=true,
    breakatwhitespace=false,
    tabsize=4,
    showstringspaces=false,
    frame=single,
    framesep=4pt,
    xleftmargin=16pt,
    xrightmargin=4pt,
    captionpos=b,
}

\lstset{style=cstyle}

% ── en-tête et pied de page ───────────────────────────────────────────────────
\pagestyle{fancy}
\fancyhf{}
\rhead{Sockets 2 --- Programmation Système}
\lhead{Chat TCP Multithreadé}
\cfoot{\thepage}

% ── document ──────────────────────────────────────────────────────────────────
\begin{document}

\begin{titlepage}
    \centering
    \vspace*{2cm}
    {\Huge\bfseries Chat TCP Multithreadé\par}
    \vspace{0.5cm}
    {\Large Programmation Système en C --- Sockets 2\par}
    \vspace{2cm}
    \rule{\linewidth}{0.4pt}\\[0.4cm]
    {\large Chat full-duplex sur socket TCP avec fils d'exécution POSIX\par}
    \rule{\linewidth}{0.4pt}
    \vfill
    {\large Ilyass Bougati\par}
    \vspace{0.3cm}
    {\large \today\par}
\end{titlepage}

\tableofcontents
\newpage

% ── Introduction ──────────────────────────────────────────────────────────────
\section{Introduction}

Ce projet est l'évolution directe du talkie-walkie (\emph{Sockets 1}). Il
lève la contrainte semi-duplex en rendant la communication \textbf{full-duplex}
: les deux interlocuteurs peuvent écrire et recevoir des messages en même temps,
sans attendre leur tour.

Pour cela, chaque côté (client et serveur) crée un fil d'exécution
(\texttt{pthread}) dédié à la \emph{réception}, tandis que le fil principal
se charge de l'\emph{envoi}. Les deux fils partagent un drapeau
\texttt{volatile int running} pour se coordonner lors de la déconnexion.

% ── Architecture ──────────────────────────────────────────────────────────────
\section{Architecture}

\subsection{Modèle de threading}

\begin{center}
\begin{tabular}{@{}lll@{}}
\toprule
Fil & Rôle & Terminaison \\
\midrule
Principal (main) & Lire stdin et envoyer (\texttt{send}) & \texttt{/quit} ou EOF \\
\texttt{recv\_thread} & Recevoir et afficher (\texttt{recv}) & Déconnexion ou \texttt{/quit} distant \\
\bottomrule
\end{tabular}
\end{center}

Les deux fils partagent la structure \texttt{thread\_arg\_t} :

\begin{lstlisting}[caption={Structure partagée entre les deux fils}]
typedef struct {
    int fd;
    char peer[NAME_SIZE];
    volatile int *running;
} thread_arg_t;
\end{lstlisting}

Le champ \texttt{running} est déclaré \texttt{volatile} pour indiquer au
compilateur qu'il peut être modifié depuis un autre fil et empêcher toute
optimisation qui le mettrait en cache dans un registre.

\subsection{Échange des pseudonymes}

Le handshake est identique à celui du talkie-walkie :

\begin{enumerate}
    \item Le serveur envoie son nom en premier.
    \item Le client le reçoit, saisit le sien et l'envoie en retour.
\end{enumerate}

Une fois les noms échangés, le fil de réception est lancé et la boucle
d'envoi démarre simultanément.

\subsection{Arrêt propre}

L'arrêt coordonné des deux fils est le point le plus délicat du programme.
Lorsqu'un fil détecte la fin de la conversation (commande \texttt{/quit} ou
déconnexion réseau), il positionne \texttt{running = 0}. Le fil principal
appelle ensuite :

\begin{lstlisting}[caption={Séquence d'arrêt}]
running = 0;
shutdown(sock_fd, SHUT_RDWR);  /* debloque le recv bloquant */
pthread_join(tid, NULL);       /* attendre la fin du fil de reception */
\end{lstlisting}

L'appel à \texttt{shutdown(SHUT\_RDWR)} est essentiel : il force le
\texttt{recv} bloquant dans \texttt{recv\_thread} à retourner immédiatement
avec une valeur nulle ou négative, ce qui permet au fil de se terminer
proprement sans attendre un prochain message réseau.

\subsection{Affichage sans collision}

Lorsqu'un message entrant arrive pendant que l'utilisateur est en train de
taper, \texttt{recv\_thread} utilise la séquence d'échappement ANSI
\texttt{\textbackslash r\textbackslash 033[K} pour revenir en début de ligne
et l'effacer avant d'afficher le message reçu. Cela évite que le texte en
cours de saisie et le message entrant se mélangent visuellement.

% ── Appels système ────────────────────────────────────────────────────────────
\section{Appels système et API utilisés}

\begin{description}
    \item[\texttt{socket} / \texttt{bind} / \texttt{listen} / \texttt{accept} / \texttt{connect}]
        Infrastructure TCP standard, identique à \emph{Sockets 1}.
    \item[\texttt{send} / \texttt{recv}]
        Envoi et réception de données sur la socket connectée.
    \item[\texttt{pthread\_create} / \texttt{pthread\_join}]
        Création et attente du fil de réception (API POSIX Threads).
    \item[\texttt{shutdown}]
        Fermeture unidirectionnelle ou bidirectionnelle d'une socket,
        utilisée ici pour débloquer un \texttt{recv} en attente.
    \item[\texttt{signal(SIGINT, ...)}]
        Gestion de Ctrl+C pour fermer proprement le descripteur d'écoute
        côté serveur.
    \item[\texttt{setsockopt(SO\_REUSEADDR)}]
        Permet de relancer le serveur immédiatement après un arrêt sans
        attendre l'expiration du délai \texttt{TIME\_WAIT}.
\end{description}

% ── Utilisation ───────────────────────────────────────────────────────────────
\section{Utilisation}

\begin{lstlisting}[language=bash, style=cstyle]
# Terminal 1 : lancer le serveur sur le port 4242
./server 4242

# Terminal 2 : connecter le client
./client 127.0.0.1 4242
\end{lstlisting}

La compilation requiert l'option \texttt{-pthread} (déjà présente dans le
\texttt{Makefile}) pour lier la bibliothèque POSIX Threads.

% ── Code source ───────────────────────────────────────────────────────────────
\section{Code source}

\subsection{\texttt{server.c}}

\begin{lstlisting}[caption={\texttt{server.c} --- serveur de chat full-duplex multithreadé}]
#include <stdio.h>
#include <stdlib.h>
#include <string.h>
#include <unistd.h>
#include <signal.h>
#include <pthread.h>
#include <arpa/inet.h>

#define BUF_SIZE 1024
#define NAME_SIZE 64

static int server_fd = -1;

typedef struct {
    int fd;
    char peer[NAME_SIZE];
    volatile int *running;
} thread_arg_t;

static void handle_sigint(int sig) {
    (void)sig;
    printf("\nClosing the server application...\n");
    if (server_fd != -1)
        close(server_fd);
    exit(0);
}

static void *recv_thread(void *arg) {
    thread_arg_t *a = arg;
    char buf[BUF_SIZE];
    ssize_t n;

    while (*a->running) {
        n = recv(a->fd, buf, sizeof(buf) - 1, 0);
        if (n <= 0) {
            if (*a->running)
                printf("\n%s disconnected. Press Enter to exit.\n", a->peer);
            *a->running = 0;
            return NULL;
        }
        buf[n] = '\0';
        if (strcmp(buf, "/quit") == 0) {
            printf("\n%s left the chat. Press Enter to exit.\n", a->peer);
            *a->running = 0;
            return NULL;
        }
        printf("\r\033[K%s: %s\n", a->peer, buf);
        fflush(stdout);
    }
    return NULL;
}

static void chat(int client_fd, const char *peer_ip) {
    char my_name[NAME_SIZE], their_name[NAME_SIZE], buf[BUF_SIZE];
    ssize_t n;
    volatile int running = 1;

    printf("Enter your username: ");
    fflush(stdout);
    if (!fgets(my_name, sizeof(my_name), stdin))
        return;
    my_name[strcspn(my_name, "\n")] = '\0';

    send(client_fd, my_name, strlen(my_name), 0);

    n = recv(client_fd, their_name, sizeof(their_name) - 1, 0);
    if (n <= 0)
        return;
    their_name[n] = '\0';

    printf("Connected with %s (%s). You can both type freely!\n\n",
           their_name, peer_ip);

    thread_arg_t arg = { .fd = client_fd, .running = &running };
    strncpy(arg.peer, their_name, NAME_SIZE - 1);
    arg.peer[NAME_SIZE - 1] = '\0';

    pthread_t tid;
    pthread_create(&tid, NULL, recv_thread, &arg);

    while (running) {
        printf("%s: ", my_name);
        fflush(stdout);
        if (!fgets(buf, sizeof(buf), stdin) || !running)
            break;
        buf[strcspn(buf, "\n")] = '\0';
        if (!running)
            break;
        send(client_fd, buf, strlen(buf), 0);
        if (strcmp(buf, "/quit") == 0) {
            running = 0;
            break;
        }
    }

    running = 0;
    shutdown(client_fd, SHUT_RDWR);
    pthread_join(tid, NULL);
    close(client_fd);
    printf("Chat ended. Waiting for a new connection...\n\n");
}

int main(int argc, char *argv[]) {
    if (argc != 2) {
        fprintf(stderr, "Usage: %s <port>\n", argv[0]);
        exit(EXIT_FAILURE);
    }

    int port = atoi(argv[1]);
    if (port <= 0 || port > 65535) {
        fprintf(stderr, "Invalid port: %s\n", argv[1]);
        exit(EXIT_FAILURE);
    }

    signal(SIGINT, handle_sigint);

    server_fd = socket(AF_INET, SOCK_STREAM, 0);
    if (server_fd < 0) {
        perror("socket");
        exit(EXIT_FAILURE);
    }

    int opt = 1;
    setsockopt(server_fd, SOL_SOCKET, SO_REUSEADDR, &opt, sizeof(opt));

    struct sockaddr_in address = {
        .sin_family      = AF_INET,
        .sin_addr.s_addr = INADDR_ANY,
        .sin_port        = htons(port)
    };

    if (bind(server_fd, (struct sockaddr *)&address, sizeof(address)) < 0) {
        perror("bind");
        exit(EXIT_FAILURE);
    }

    if (listen(server_fd, 1) < 0) {
        perror("listen");
        exit(EXIT_FAILURE);
    }

    printf("Server listening on port %d...\n\n", port);

    socklen_t addr_len = sizeof(address);
    while (1) {
        int client_fd = accept(server_fd,
                               (struct sockaddr *)&address, &addr_len);
        if (client_fd < 0) {
            perror("accept");
            continue;
        }
        printf("Client connected from %s\n", inet_ntoa(address.sin_addr));
        chat(client_fd, inet_ntoa(address.sin_addr));
    }

    close(server_fd);
    return 0;
}
\end{lstlisting}

\subsection{\texttt{client.c}}

\begin{lstlisting}[caption={\texttt{client.c} --- client de chat full-duplex multithreadé}]
#include <stdio.h>
#include <stdlib.h>
#include <string.h>
#include <unistd.h>
#include <pthread.h>
#include <arpa/inet.h>

#define BUF_SIZE 1024
#define NAME_SIZE 64

typedef struct {
    int fd;
    char peer[NAME_SIZE];
    volatile int *running;
} thread_arg_t;

static void *recv_thread(void *arg) {
    thread_arg_t *a = arg;
    char buf[BUF_SIZE];
    ssize_t n;

    while (*a->running) {
        n = recv(a->fd, buf, sizeof(buf) - 1, 0);
        if (n <= 0) {
            if (*a->running)
                printf("\n%s disconnected. Press Enter to exit.\n", a->peer);
            *a->running = 0;
            return NULL;
        }
        buf[n] = '\0';
        if (strcmp(buf, "/quit") == 0) {
            printf("\n%s left the chat. Press Enter to exit.\n", a->peer);
            *a->running = 0;
            return NULL;
        }
        printf("\r\033[K%s: %s\n", a->peer, buf);
        fflush(stdout);
    }
    return NULL;
}

static void chat(int sock_fd) {
    char my_name[NAME_SIZE], their_name[NAME_SIZE], buf[BUF_SIZE];
    ssize_t n;
    volatile int running = 1;

    n = recv(sock_fd, their_name, sizeof(their_name) - 1, 0);
    if (n <= 0) {
        fprintf(stderr, "Server closed connection.\n");
        return;
    }
    their_name[n] = '\0';

    printf("Enter your username: ");
    fflush(stdout);
    if (!fgets(my_name, sizeof(my_name), stdin))
        return;
    my_name[strcspn(my_name, "\n")] = '\0';

    send(sock_fd, my_name, strlen(my_name), 0);

    printf("Connected with %s. You can both type freely!\n\n", their_name);

    thread_arg_t arg = { .fd = sock_fd, .running = &running };
    strncpy(arg.peer, their_name, NAME_SIZE - 1);
    arg.peer[NAME_SIZE - 1] = '\0';

    pthread_t tid;
    pthread_create(&tid, NULL, recv_thread, &arg);

    while (running) {
        printf("%s: ", my_name);
        fflush(stdout);
        if (!fgets(buf, sizeof(buf), stdin) || !running)
            break;
        buf[strcspn(buf, "\n")] = '\0';
        if (!running)
            break;
        send(sock_fd, buf, strlen(buf), 0);
        if (strcmp(buf, "/quit") == 0) {
            running = 0;
            break;
        }
    }

    running = 0;
    shutdown(sock_fd, SHUT_RDWR);
    pthread_join(tid, NULL);
}

int main(int argc, char *argv[]) {
    if (argc != 3) {
        fprintf(stderr, "Usage: %s <host> <port>\n", argv[0]);
        exit(EXIT_FAILURE);
    }

    int port = atoi(argv[2]);
    if (port <= 0 || port > 65535) {
        fprintf(stderr, "Invalid port: %s\n", argv[2]);
        exit(EXIT_FAILURE);
    }

    int sock_fd = socket(AF_INET, SOCK_STREAM, 0);
    if (sock_fd < 0) {
        perror("socket");
        exit(EXIT_FAILURE);
    }

    struct sockaddr_in addr = {
        .sin_family = AF_INET,
        .sin_port   = htons(port)
    };

    if (inet_pton(AF_INET, argv[1], &addr.sin_addr) <= 0) {
        fprintf(stderr, "Invalid address: %s\n", argv[1]);
        exit(EXIT_FAILURE);
    }

    if (connect(sock_fd, (struct sockaddr *)&addr, sizeof(addr)) < 0) {
        perror("connect");
        exit(EXIT_FAILURE);
    }

    chat(sock_fd);
    close(sock_fd);
    return 0;
}
\end{lstlisting}

\subsection{Makefile}

\begin{lstlisting}[language=make, caption={Makefile --- compilation avec support POSIX Threads}]
CC = gcc
CFLAGS = -Wall -Wextra -pthread

all: server client

server: server.c
	$(CC) $(CFLAGS) -o $@ $<

client: client.c
	$(CC) $(CFLAGS) -o $@ $<

clean:
	rm -f server client

.PHONY: all clean
\end{lstlisting}

% ── Comparaison Sockets 1 / Sockets 2 ────────────────────────────────────────
\section{Comparaison avec Sockets 1}

\begin{center}
\begin{tabular}{@{}lll@{}}
\toprule
Aspect & Sockets 1 (talkie-walkie) & Sockets 2 (multithreadé) \\
\midrule
Mode de communication & Semi-duplex (tours alternés) & Full-duplex (simultané) \\
Threads & Aucun & 1 fil de réception par côté \\
Blocage & \texttt{recv} bloque l'envoi & Les deux opèrent en parallèle \\
Arrêt & Retour de boucle simple & \texttt{shutdown} + \texttt{pthread\_join} \\
Complexité & Faible & Moyenne \\
\bottomrule
\end{tabular}
\end{center}

% ── Bilan ─────────────────────────────────────────────────────────────────────
\section{Bilan}

Ce projet démontre comment les fils d'exécution POSIX permettent de lever la
contrainte bloquante inhérente aux appels \texttt{recv}, transformant un chat
en mode talkie-walkie en une véritable messagerie en temps réel. Les points
clés à retenir sont l'utilisation de \texttt{volatile} pour la variable
partagée entre fils, l'importance de \texttt{shutdown} pour débloquer un
\texttt{recv} en attente, et la nécessité d'attendre la fin du fil de
réception via \texttt{pthread\_join} avant de fermer le descripteur de socket.

\end{document}
